% ============================================================
% Project: THE PHILOSOPHY OF VIABILITY
% Book: Book I — Ontology of Viability
% Book ID: VS.B01
% Chapter: Chapter 07 — Formal and Logical Reconstruction
% File: 4_ASSETS/tables/VS.B01.P01.C07.S03.SS04.P001.TL01.formal_domains_and_types.tex
% Asset: TL01 — Formal Domains and Types
% Purpose: Summary of formal expression types and their
% interpretive and ontological constraints.
% ============================================================

\begin{table}[htbp]
\centering
\small
\caption{Formal Expression Types and Their Roles}
\label{tab:ch07_formal_types}
\renewcommand{\arraystretch}{1.25}
\begin{tabularx}{\textwidth}{|p{0.19\textwidth}|p{0.27\textwidth}|X|}
\hline
\textbf{Formal Type} &
\textbf{Formal Role} &
\textbf{Interpretive and Ontological Constraint}
\\
\hline
Object variable &
Ranges over a specified formal domain. &
Its range must be stated; it does not establish the complete inventory of existents.
\\
\hline
Constant &
Designates an element under an interpretation. &
Designation within a formal model does not independently establish ontological status.
\\
\hline
Predicate &
Expresses a condition attributed to one or more arguments. &
Formal predication does not establish whether the attributed condition is essential or identity-constitutive.
\\
\hline
Relation &
Connects arguments of specified types and arity. &
Its formal properties must not be transferred automatically to the represented ontological relation.
\\
\hline
Function &
Assigns a unique output to each admissible input. &
Totality and uniqueness must be justified; otherwise a partial function or relation may be required.
\\
\hline
Operator &
Acts upon formal expressions according to specified rules. &
Formal operation must not be mistaken for an ontological agent or process.
\\
\hline
Transformation &
Represents a transition between conditions or configurations. &
Must specify its inputs, outputs, domain, and whether it is functional or relational.
\\
\hline
Proposition &
Expresses content that may be asserted, denied, or assessed. &
Must be distinguished from the objects and expressions about which it makes claims.
\\
\hline
\end{tabularx}
\end{table}

% ============================================================
% END OF ASSET TL01 — CHAPTER 07, SECTION 03
% ============================================================